\chapter{S5}
\label{ch:s5}

\section{Why S5 Exists}

S4 showed that structured state space layers could model long sequences effectively, but
it came with significant machinery: normal-plus-low-rank matrices, complex kernel
generation, Cauchy computations, and a fairly intricate implementation. S5, short for
Simplified State Space Layers, keeps the core SSM idea while simplifying the architecture
and computation \citep{smith2023s5}.

At a high level, S5 makes three important moves:

\begin{itemize}
  \item use diagonal state matrices rather than S4's more elaborate NPLR form;
  \item use multi-input multi-output SSMs rather than many independent SISO filters;
  \item use parallel scan to evaluate the recurrence efficiently.
\end{itemize}

This makes S5 a particularly good model to understand after S4. It preserves the
state-space perspective, but many pieces become easier to reason about.

\section{The S5 Recurrence}

The discrete SSM recurrence is still the starting point:
\[
  \vx_t = \bar{\mA}\vx_{t-1} + \bar{\mB}\vu_t,
  \qquad
  \vy_t = \mC\vx_t + \mD\vu_t.
\]
S5 uses a diagonal state transition:
\[
  \bar{\mA} = \diag(\bar{\lambda}_1,\ldots,\bar{\lambda}_N).
\]
Then each state coordinate evolves independently before the readout:
\[
  x_t[j] = \bar{\lambda}_j x_{t-1}[j] + (\bar{\mB}\vu_t)[j].
\]
The coordinates interact through $\bar{\mB}$ and $\mC$, not through a dense transition
matrix. This is much simpler than a generic dense $\bar{\mA}$.

The kernel is still
\[
  \mK_i = \mC\bar{\mA}^i\bar{\mB}.
\]
With diagonal $\bar{\mA}$,
\[
  \bar{\mA}^i
  =
  \diag(\bar{\lambda}_1^i,\ldots,\bar{\lambda}_N^i),
\]
so the kernel is a mixture of independent exponential modes.

\section{Diagonal Does Not Mean Trivial}

It is tempting to think a diagonal state matrix is too simple. But diagonal dynamics can
still be expressive when combined with input and output mixing.

The input matrix $\bar{\mB}$ maps the input vector into all state coordinates:
\[
  \bar{\mB}\vu_t \in \C^N.
\]
The output matrix $\mC$ mixes all state coordinates back into output channels:
\[
  \mC\vx_t \in \R^{C_{\mathrm{out}}}.
\]
So each output channel can depend on many state modes, and each input channel can excite
many state modes. The state modes themselves are independent, but the layer as a whole
can still mix channels richly.

This is the same principle that makes diagonalization useful in linear systems. In the
right basis, dynamics are independent modes; input and output projections determine how
those modes are used.

\section{Complex Diagonal Modes}

S5 commonly uses complex diagonal modes. Each eigenvalue-like parameter can be written as
\[
  \lambda_j = \alpha_j + i\omega_j.
\]
In continuous time, stability requires
\[
  \RePart(\lambda_j)=\alpha_j < 0.
\]
After discretization with step size $\Delta$, the discrete eigenvalue is
\[
  \bar{\lambda}_j = e^{\Delta \lambda_j}.
\]
Its magnitude is
\[
  |\bar{\lambda}_j| = e^{\Delta \alpha_j},
\]
so negative real part gives a magnitude below $1$. The imaginary part $\omega_j$ controls
oscillation.

Thus, S5's diagonal state can represent a bank of decaying oscillatory filters:
\[
  \bar{\lambda}_j^t
  =
  e^{\Delta \alpha_j t} e^{i\Delta\omega_j t}.
\]
This is a powerful basis for temporal modeling.

\section{MIMO State Space Layers}

One of S5's main changes from S4 is the move toward multi-input multi-output SSMs. Instead
of treating channels as many independent SISO filters, S5 allows a single SSM layer to map
multiple input channels to multiple output channels:
\[
  \vu_t \in \R^{H},
  \qquad
  \vy_t \in \R^{H},
  \qquad
  \vx_t \in \C^N.
\]
Here $H$ is model width and $N$ is state size.

The parameter shapes are
\[
  \bar{\mB}\in\C^{N\times H},
  \qquad
  \mC\in\C^{H\times N}.
\]
The resulting kernel has shape
\[
  \mK_i\in\R^{H\times H}
\]
after taking the real-valued output convention. So at each lag $i$, the layer can map all
input channels to all output channels.

This MIMO view is appealing because it makes the SSM layer feel more like a standard
sequence block over model-width features rather than a collection of independent scalar
filters plus separate mixing.

\section{Discretization And Learned Timescales}

Like S4, S5 can start from continuous-time parameters and discretize them. For diagonal
$\mA$, discretization is especially simple:
\[
  \bar{\lambda}_j = e^{\Delta_j \lambda_j}.
\]
The step size $\Delta_j$ may be learned or parameterized. This gives each mode its own
effective timescale.

The input matrix also needs discretization. For a scalar diagonal mode,
\[
  \bar{B}_j
  =
  \left(\int_0^{\Delta_j} e^{\tau\lambda_j}d\tau\right)B_j
  =
  \frac{e^{\Delta_j\lambda_j}-1}{\lambda_j}B_j
\]
when $\lambda_j\neq 0$. The exact implementation may use numerically stable variants,
but the idea is the same as Chapter~\ref{ch:state-space-models}: continuous dynamics are
converted into discrete sequence updates.

\section{Parallel Scan}

S5 emphasizes parallel scan as a way to evaluate the recurrence efficiently. To understand
scan, start with a recurrence of the form
\[
  \vx_t = \mA_t\vx_{t-1}+\vb_t.
\]
For S5 with fixed diagonal $\bar{\mA}$,
\[
  \mA_t=\bar{\mA},
  \qquad
  \vb_t=\bar{\mB}\vu_t.
\]

Each time step defines an affine map:
\[
  F_t(\vx)=\mA_t\vx+\vb_t.
\]
Composing two adjacent updates gives
\[
  F_2(F_1(\vx))
  =
  \mA_2\mA_1\vx + \mA_2\vb_1+\vb_2.
\]
So the composition of two affine maps is another affine map:
\[
  (\mA_2,\vb_2)\circ(\mA_1,\vb_1)
  =
  (\mA_2\mA_1,\; \mA_2\vb_1+\vb_2).
\]

This composition is associative. Associativity means we can compute all prefix
compositions with a parallel scan algorithm. That gives the same result as stepping
through the recurrence, but with much more parallelism on suitable hardware.

This matters because it gives S5 another path to efficient training. Instead of explicitly
forming a long convolution kernel, one can evaluate the recurrence through an associative
parallel prefix operation.

\section{Scan Versus Convolution}

For a time-invariant linear SSM, recurrent mode and convolution mode are mathematically
equivalent. S4 leans heavily on generating and applying the convolution kernel. S5's
parallel scan view emphasizes efficient recurrence evaluation.

The distinction is practical:

\[
\begin{array}{lll}
\text{Convolution mode} &:& \text{construct or apply } \mK_i=\mC\bar{\mA}^i\bar{\mB}, \\
\text{Scan mode} &:& \text{parallelize the recurrence updates directly}.
\end{array}
\]

Both exploit structure. Convolution exploits time invariance and fast filtering. Scan
exploits associativity of recurrence composition. Mamba will also rely heavily on scan,
because its input-dependent dynamics break the fixed convolution kernel.

\section{Initialization From S4 And HiPPO}

S5 simplifies S4, but it does not discard the long-memory initialization story. It can use
initializations derived from S4/HiPPO-style dynamics. The goal is still to start with a
useful spectrum of modes rather than random short-memory behavior.

Conceptually:

\[
  \text{HiPPO/S4 initialization}
  \rightarrow
  \text{diagonalized modes}
  \rightarrow
  \text{S5 diagonal SSM}.
\]

The practical result is that S5 begins with modes covering a range of timescales. Training
then adapts those modes, along with $\bar{\mB}$, $\mC$, step sizes, and surrounding neural
network parameters.

\section{The S5 Block}

A simplified S5 block has the same broad structure as other residual sequence blocks:
\[
  \vz
  \rightarrow
  \text{normalization}
  \rightarrow
  \text{linear projection}
  \rightarrow
  \text{S5 temporal mixing}
  \rightarrow
  \text{nonlinearity or gate}
  \rightarrow
  \text{linear projection}
  \rightarrow
  \text{residual addition}.
\]
Implementation details vary, but this is the right conceptual structure.

The S5 temporal mixing step is the diagonal MIMO SSM. The rest of the block supplies
feature mixing, nonlinearity, regularization, and stable deep optimization.

\section{State Size And Model Width}

S5 makes the distinction between state size and model width especially important.

\[
\begin{array}{ll}
H & \text{number of feature channels passed between blocks}, \\
N & \text{number of SSM state modes}.
\end{array}
\]

Increasing $H$ gives the network more representation channels. Increasing $N$ gives the
SSM more temporal modes. In a diagonal SSM, each state coordinate is one mode, so $N$ is
directly tied to the richness of the temporal filter bank.

For neural time series, this distinction is practically important. A model might need
enough width to represent neural features and enough state size to capture dynamics
across speech-relevant timescales. These are related but not identical forms of capacity.

\section{S5 Compared With S4}

The cleanest comparison is:

\[
\begin{array}{lll}
\text{S4} &:& \text{NPLR structure, SISO emphasis, specialized kernel computation}, \\
\text{S5} &:& \text{diagonal structure, MIMO emphasis, parallel scan}.
\end{array}
\]

S4 is closer to the full structured HiPPO matrix story. S5 is simpler and often easier to
implement and adapt. The simplification is not free: diagonal dynamics are a stronger
restriction than NPLR dynamics. But the MIMO formulation and learnable projections can
recover substantial expressivity in practice.

For many applied projects, S5 is attractive because it provides the core SSM advantages
without requiring the full S4 implementation stack.

\section{Why S5 Is Interesting For Neural Data}

Neural time series are not token sequences in the language-modeling sense. They are
continuous, noisy, dynamical signals. SSMs are a natural inductive bias because they model
temporal evolution through latent state.

S5 is especially appealing for this setting because:

\begin{itemize}
  \item it is causal by construction when run left to right;
  \item it can maintain fixed-size recurrent state for online decoding;
  \item diagonal modes provide interpretable timescale and oscillation structure;
  \item MIMO input/output maps can mix many neural channels;
  \item parallel scan makes training more practical than naive recurrence.
\end{itemize}

This does not guarantee good downstream decoding. The local project notes already show
that SSL objective design, nuisance structure, normalization, and evaluation matter a
great deal. But architecturally, S5 is a reasonable backbone to test for causal neural
sequence modeling.

\section{Limitations}

S5 also has limitations.

\paragraph{Fixed dynamics.}
Like S4, a standard S5 layer is linear time invariant inside the temporal core. The same
dynamics apply regardless of input content. This is efficient, but less flexible than
attention's content-dependent routing.

\paragraph{Diagonal state transition.}
Diagonal dynamics are simple and efficient, but they restrict interactions between state
modes. Mode interaction must happen through input/output projections and nonlinear
stacking across layers.

\paragraph{Objective dependence.}
A good SSM backbone does not by itself solve self-supervised learning. If an objective can
be solved through local recording-context shortcuts or by predicting a trivial normalized
mean, S5 may learn unhelpful representations.

\section{A Reading Strategy For The S5 Paper}

When reading \citet{smith2023s5}, separate the ideas into four layers:

\begin{enumerate}
  \item The simplification layer: replace S4's harder structured matrices with diagonal
  state dynamics.
  \item The channel layer: use MIMO SSMs for richer input-output mixing.
  \item The algorithm layer: use parallel scan to evaluate recurrences efficiently.
  \item The architecture layer: place S5 inside a deep residual sequence model.
\end{enumerate}

If an equation feels hard, ask which of these layers it belongs to. That usually makes the
paper easier to navigate.

\section{What To Remember}

\begin{itemize}
  \item S5 is a simplified state space sequence layer built after S4.
  \item It uses diagonal state transitions, often with complex-valued modes.
  \item Diagonal modes represent decays and oscillations.
  \item MIMO $\bar{\mB}$ and $\mC$ matrices let the layer mix many input and output
  channels.
  \item Parallel scan evaluates recurrence updates efficiently using associativity.
  \item S5 preserves the recurrent/convolutional SSM spirit while reducing S4's
  implementation complexity.
  \item For neural sequence modeling, S5 is attractive because it is causal, efficient,
  and organized around latent dynamics.
\end{itemize}

\section{Exercises}

\begin{enumerate}
  \item For diagonal $\bar{\mA}=\diag(\lambda_1,\lambda_2)$, write
  $\bar{\mA}^i$ explicitly.
  \item Explain why diagonal state dynamics can still produce rich outputs when paired
  with dense $\bar{\mB}$ and $\mC$ matrices.
  \item Show that composing affine recurrence updates
  $(\mA_2,\vb_2)\circ(\mA_1,\vb_1)$ gives
  $(\mA_2\mA_1,\mA_2\vb_1+\vb_2)$.
  \item In your own words, compare convolution mode and scan mode.
  \item Why might fixed LTI dynamics be limiting for tasks that require content-dependent
  selection?
\end{enumerate}

\section{Prepares You For}

Mamba modifies the SSM idea by making the dynamics selective, meaning some SSM parameters
depend on the current input. This gives up the fixed convolution kernel but keeps an
efficient scan-based recurrent computation.
